\begin{statementbox}
	\textbf{\textcolor{CStatement}{条件}}
	\begin{description}[style=nextline, leftmargin=2.8em, labelsep=0.5em, itemsep=0.35em]
		\item[\textbf{\textcolor{CStatement}{条件 1（任意集合）：}}]
			设 $A,B,C$ 为全集 $U$ 中的任意子集（或任意集合），无其他限制。
	\end{description}
	\textbf{\textcolor{CStatement}{结论}}
	\begin{description}[style=nextline, leftmargin=2.8em, labelsep=0.5em, itemsep=0.45em]
		\item[\textbf{\textcolor{CStatement}{结论一（交换律）：}}]
			\textbf{\textcolor{CStatement}{直观描述：}}
			并运算与交运算的结果与参与运算的顺序无关。
			\[
				A\cup B = B\cup A,
				\qquad
				A\cap B = B\cap A
			\]
		\item[\textbf{\textcolor{CStatement}{结论二（结合律）：}}]
			\textbf{\textcolor{CStatement}{直观描述：}}
			连续的同种运算可任意添加或移动括号。
			\[
				(A\cup B)\cup C = A\cup(B\cup C),
				\qquad
				(A\cap B)\cap C = A\cap(B\cap C)
			\]
		\item[\textbf{\textcolor{CStatement}{结论三（交对并分配）：}}]
			\textbf{\textcolor{CStatement}{直观描述：}}
			先并再交等于先交再并（交对并“穿入”）。
			\[
				A\cap(B\cup C) = (A\cap B)\cup(A\cap C)
			\]
		\item[\textbf{\textcolor{CStatement}{结论四（并对交分配）：}}]
			\textbf{\textcolor{CStatement}{直观描述：}}
			先交再并等于先并再交（并对交“穿入”）。
			\[
				A\cup(B\cap C) = (A\cup B)\cap(A\cup C)
			\]
	\end{description}
	\textbf{\textcolor{CStatement}{推广结论（有限推广）}}
	\textbf{\textcolor{CStatement}{直观描述：}}
	上述运算律对任意有限个集合同样适用，可结合数学归纳法使用。 例如，对 $n$ 个集合的并 $\bigcup_{i=1}^{n}A_i$，结合律保证任一加括号方式不改变结果；分配律可推广为
	\[
		A\cap\left(\bigcup_{i=1}^{n}B_i\right) = \bigcup_{i=1}^{n}(A\cap B_i),\quad A\cup\left(\bigcap_{i=1}^{n}B_i\right) = \bigcap_{i=1}^{n}(A\cup B_i).
	\]
	\textbf{\textcolor{CStatement}{使用提醒：}}
	空集 $\varnothing$ 也满足所有运算律，运用时无需单独讨论。
\end{statementbox}
